\subsection{\texorpdfstring{第~\ref{sec:integration-on-manifolds}~节}{流形上的积分}}\label{subsec:solutions-integration-on-manifolds}

\begin{description}[font=\normalfont,leftmargin=2em,labelwidth=1.5em]
\RGsolution{ex:integral-on-complement-of-torus-in-sphere}{integral-on-complement-of-torus-in-sphere}
$T^2$ 将 $S^3$ 分成两个连通分支
\[
  M_1:=\{(x,y,z,w)\in S^3\mid x^2+y^2<1\}
\]
和
\[
  M_2:=\{(x,y,z,w)\in S^3\mid x^2+y^2>1\}.
\]
考虑参数化 $\varphi_i:U_i\to M_i$ ($i=1,2$), 在 $U_1:=(0,1)\times(0,2\pi)\times(0,2\pi)$ 和 $U_2:=(1,\sqrt2)\times(0,2\pi)\times(0,2\pi)$ 上由
\[
  \varphi_i(r,v,v)=\left(r\cos v,r\sin v,\sqrt{2-r^2}\cos v,\sqrt{2-r^2}\sin v\right)
\]
定义. 则
\begin{align*}
  \varphi_i^*\mathrm{d}x&=\cos v\,\mathrm{d}r-r\sin v\,\mathrm{d}v\\
  \varphi_i^*\mathrm{d}y&=\sin v\,\mathrm{d}r+r\cos v\,\mathrm{d}v\\
  \varphi_i^*\mathrm{d}z&=-\frac{r}{\sqrt{2-r^2}}\cos v\,\mathrm{d}r-\sqrt{2-r^2}\sin v\,\mathrm{d}v\\
  \varphi_i^*\mathrm{d}w&=-\frac{r}{\sqrt{2-r^2}}\sin v\,\mathrm{d}r+\sqrt{2-r^2}\cos v\,\mathrm{d}v,
\end{align*}
因而 $\varphi_i^*\omega=(r(2-r^2)\cos^2 v-r^3\cos^2 v)\,\mathrm{d}r\wedge\mathrm{d}v\wedge\mathrm{d}v$. 由于 $M_1\setminus\varphi_1(U_1)$ 和 $M_2\setminus\varphi_2(U_2)$ 的测度为零, 有
\begin{align*}
  \int_{M_1}\omega&=\pm\int_{U_1}\varphi_1^*\omega\\
  &=\pm\int_0^1\int_0^{2\pi}\int_0^{2\pi}
       (r(2-r^2)\cos^2 v-r^3\cos^2 v)\,\mathrm{d}v\,\mathrm{d}v\,\mathrm{d}r=\pm\pi^2,\\
  \int_{M_2}\omega&=\pm\int_{U_2}\varphi_2^*\omega\\
  &=\pm\int_1^{\sqrt{2}}\int_0^{2\pi}\int_0^{2\pi}
       (r(2-r^2)\cos^2 v-r^3\cos^2 v)\,\mathrm{d}v\,\mathrm{d}v\,\mathrm{d}r=\mp\pi^2,
\end{align*}
其中正负号取决于定向的选取.

\RGsolution{ex:orientation-preserving-pullback-integral}{orientation-preserving-pullback-integral}
考虑 $M$ 上由保持定向的参数化构成的一个图册 $\{(U_\alpha,\varphi_\alpha)\}$. 由于 $f$ 是保持定向的微分同胚, $\{(U_\alpha,f\circ\varphi_\alpha)\}$ 是 $N$ 上由保持定向的参数化构成的图册. 设 $\{\rho_i\}_{i\in I}$ 是从属于 $N$ 的覆盖 $\{W_\alpha\}$ 的单位分解, 其中 $W_\alpha:=(f\circ\varphi_\alpha)(U_\alpha)$. 由于 $\omega$ 有紧支集, 可以假设 $I$ 是有限的. 因此, 对 $\omega_i:=\rho_i\omega$, 有
\[
  \int_N\omega:=\sum_{i\in I}\int_N\omega_i
    =\sum_{i\in I}\int_{U_{\alpha_i}}(f\circ\varphi_{\alpha_i})^*\omega_i.
\]
现在,
\begin{align*}
  (f\circ\varphi_{\alpha_i})^*\omega_i
  &=(f\circ\varphi_{\alpha_i})^*(\rho_i\omega)
   =((f\circ\varphi_{\alpha_i})^*\rho_i)((f\circ\varphi_{\alpha_i})^*\omega)\\
  &=(\rho_i\circ f\circ\varphi_{\alpha_i})(\varphi_{\alpha_i}^*(f^*\omega))
   =\varphi_{\alpha_i}^*((\rho_i\circ f)f^*\omega).
\end{align*}
此外, $\{\rho_i\circ f\}_{i\in I}$ 显然是从属于 $M$ 的覆盖 $\{W_\alpha\}$ 的单位分解, 其中 $W_\alpha:=\varphi_\alpha(U_\alpha)$. 因此,
\begin{align*}
  \int_N\omega
  &=\sum_{i\in I}\int_{U_{\alpha_i}}(f\circ\varphi_{\alpha_i})^*\omega_i
   =\sum_{i\in I}\int_{U_{\alpha_i}}\varphi_{\alpha_i}^*((\rho_i\circ f)f^*\omega)\\
  &=\sum_{i\in I}\int_{U_{\alpha_i}}\varphi_{\alpha_i}^*(f^*\omega)_i
   =\int_M f^*\omega.
\end{align*}
\end{description}
